\section{Contiki-NG and COOJA Simulator}
\subsection{Contiki}
% TODO: cite ContikiOS
Contiki-NG is a lightweight operating system for resource-constrained
devices in WSNs and IoT.

which was built on the Contiki.

The Swedish Institute of Computer Science (SICS) released the first
version of Contiki in 2003 % TODO: cite this
Contiki was created by Adam Dunkels in 2002, which supports various
protocol for IoT devices
and have been support from various research instutitions
Later on, ContikiMAC was introduce as the sole duty cycle MAC layer for RPL.

% TODO: paraphrase this paragraph
However, recently a new Contiki cross-platform has been released on
Nov 2017 named Contiki-NG 4.0, which licensed under the 3-clause BSD
license. Contiki-NG concentrates on the Next-Generation IoT devices
such as ARM Cortex M3 and other 32-bit MCUs and on reliable and
secure standard protocols like IPv6/6LoWPAN, 6TiSCH, RPL, and CoAP
\cite{ContikiNG}.

In Contiki-NG, RPL implementation was rewritten as RPL-lite while the
original version of RPL in Contiki is kept as RPL-classic. RPL-lite
removes certain features RPL based on year of experience, keeping the
supports one single DAG at a time, one single “RPL instance”, and
only the non-storing mode of operation.

Addionally, ContikiMAC was removed in favor of TSCH implemetation,
thus making TSCH the default for duty cycle MAC implementation.

\subsection{COOJA Simulator for Contiki-NG}
The Cooja simulator is widely used among RPL OF researches % TODO:
% reference here

Cooja is an elastic Java-based simulator dedicated for WSNs based on
Contiki Operating System.
Cooja is capable of simulating different node types, including
simulated hardward (consisting of clock cycle and MCU) and software
based simulation (which skip the internal working of an MCU).

In this regard, Cooja is a tool for real hardware simulation.

depends on the mote type, provides three mote types out-of-the-box:
Sky, Z1 and Cooja where Sky and Z1 are emulated hardware for TelosB
and Zolertia Z1 and Cooja is the software simulation.
Table~\ref{tab:mote-specs} summarises the board specification.

% TODO: double checkthe MCU and Transceiver
\begin{table}[h]
    \caption{Supported mote hardware specifications in COOJA simulator}
    \label{tab:mote-specs}
    \begin{center}
        \begin{tabular}{lll}
            \hline
            \textbf{Specification} & \textbf{Zolertia Z1} & \textbf{TelosB} \\
            \hline
            Microcontroller        & MSP430F2617 (2nd gen)     & Texas
            Instruments MSP430 \\
            CPU                    & 16-bit RISC @ 16\,MHz     & 8\,MHz \\
            RAM                    & 8\,KB                     & 10\,KB \\
            Flash memory           & 92\,KB                    & 48\,KB \\
            Transceiver            & CC2420, IEEE 802.15.4     & Chipcon,
            IEEE 802.15.4 \\
            Frequency band         & 2.4\,GHz                  & 2.4\,GHz \\
            Data rate              & 250\,kbps                 & 250\,kbps \\
            \hline
        \end{tabular}
    \end{center}
\end{table}

Thus to maximise the computing power for machine learning model, Z1
is the most suitable mote for our need.

% TODO: paraphrase this
This platform is a virtual platform, used by the Cooja platform to
run Contiki-NG as ‘Cooja motes’. It compiles Contiki-NG as a native
process, and connects directly all hardware accesses to the Cooja
simulator. Unlike emulations, simulation with Cooja motes are not
perfectly time-accurate, but the timing is still good enough to run
the MAC layers CSMA and TSCH
